\section{Discretization}\label{sec:discretization}
% \alexandra{This description needs a clear explanation that there are two parts to the discretization: types and then functions. If I understand correctly the text before the next subsection is describing what is in Figure 6. This is completely unclear now... As far as I can see Fig.6 is not mentioned anywhere? }
%
\begin{figure}[t]
	\centering
	\footnotesize
	
	% Basic discretization rules
	\begin{tabular}{l||l}
		$e \colon \tau $ &  $\discretize{e : \tau}$ \\
		\hline
		\hline
		$ {x : \tau}$ & $x$\\[2ex]
		$  {c : \float[B;V]}$ &
		$  \begin{cases}
			c & \text{if } B = \top, \\
			\finconst{k}{n+1} \text{ such that } c \in \intervals{B}_k & \text{if } B \neq \top.
		\end{cases}$
		\\[3ex]
		$ {\letkw\; x = e_1 : \tau_1\; \inkw\; e_2 : \tau_2 : \tau_2}$ &
		$\letkw\; x = \discretize{e_1 : \tau_1}\; \inkw\; \discretize{e_2 : \tau_2}$
		\\[2ex]
		${(e_1 : \float[B_1;V_1] < e_2 : \float[B_2;V_2]) : \bool} $&
		$\discretize{e_1 : \float[B_1;V_1]} < \discretize{e_2 : \float[B_2;V_2]}$
		\\[2ex]
		$ {(e_1 : \float[B_1;V_1] \leq e_2 : \float[B_2;V_2]) : \bool}$ &
		$\discretize{e_1 : \float[B_1;V_1]} \leq \discretize{e_2 : \float[B_2;V_2]}$
		\\[2ex]
		$ {(\ifkw\; e_1 : \bool\; \thenkw\; e_2 : \tau\; \elsekw\; e_3 : \tau) : \tau}$ &
		$\ifkw\; \discretize{e_1 : \bool}\; \thenkw\; \discretize{e_2 : \tau} \elsekw\; \discretize{e_3 : \tau}$
		\\[2ex]
		$ {\uniform(e_1, e_2 ): \float[B;V]} $&\\ 
		\qquad $ \text{if } e_1, e_2 \in \mathbb{R}, e_1 > e_2 $ & $\texttt{Error: Uniform $e_1 > e_2$} $   \\[2ex]
		\qquad      $B=\top$ & $\uniform(\discretize{e_1 : \float[B_1;V_1]}, \discretize{e_2 : \float[B_2;V_2]}) $\\[2ex]
		\qquad       $\text{if } B \neq \top, 
		B = \{\sim_1\!\!b_1, \ldots, \sim_n\!\!b_n\},  \text{ and }
		e_1, e_2 \in \mathbb{R} $ 
		&            
		$\discrete(p_0, \ldots, p_n)$ \\[1.2ex]
		&$\text{ where } p_k = \text{CDF}(b_{k+1}) - \text{CDF}(b_k)$
		\\[2ex]
		\qquad   $\text{if } B \neq \top, 
		V_1 = \{v_1, \ldots, v_n\}, 
		V_2 = \{w_1, \ldots, w_m\}
		$      &
		$\letkw \; v_1' = \discretize{e_1 : \float[B_1; V_1]} \; \inkw $\\
		&$\begin{array}{@{\!\!\!\!\!\!\!}ll}
			&\letkw \; v_2' = \discretize{e_2 : \float[B_2; V_2]} \; \inkw \\
			&\ifkw\; v_1' == 0\; \&\& \; v_2' == 0\; \thenkw \\
			&\quad \discretize{\uniform(v_1, w_1)} \\
			&\vdots \\
			&\elsekw\; \ifkw\; v_1' == i\; \&\& \; v_2' == j\; \thenkw \\
			&\quad \discretize{\uniform(v_{i+1}, w_{j+1})} \\
			&\vdots \\
			&\elsekw \\
			&\quad \discretize{\uniform(v_n, w_m)}
		\end{array}
		$
	\end{tabular}
	\caption{Discretization function for \textsc{Slice}. If $e_1$ and $e_2$ are variables in \texttt{discretize($\uniform(e_1, e_2)$)}, the let bindings for v1' and v2' are in fact unnecessary; the nested \texttt{if} statements can directly use $e_1$ and $e_2$ in their guards.}
	\label{fig:discretization-code}
\end{figure}
%
In this section, we describe our type-directed discretization algorithm alongside illustrative examples. The setting is the following: Given a typed expression $e : \tau$ compute --- guided by the type annotations --- a (partially) \emph{discretized} expression $\discretize{e : \tau}$ such that $e$ and $\discretize{e : \tau}$ are \emph{semantically equivalent}, i.e., they produce the same distribution.

$\discretize{e : \tau}$ is defined recursively on $e : \tau$ as shown in \Cref{fig:discretization-code}, where we focus on continuous uniform distributions. For other distributions, this procedure is completely analogous. Before we go into the details, let us get some intuition on how we use the type annotation to discretize expressions. The type system is designed in such a way that 
%
\begin{center}
	Every typed sub-expression of the form ${\uniform(e_1, e_2 ): \float[B;V]}$ with $B=\{as\}$
\end{center}

